% abstract | Status: done | Original lines: 214-235
% REWRITTEN: synthesis count corrected (three sections + one subsection, Section sec:synth4); glossary sentence added.
%!TEX root = ../UNSORTED.tex
\begin{center}
\begin{minipage}{0.88\textwidth}
\centering{\bfseries\small Abstract}
\par\vspace{0.4em}
\small\raggedright\setlength{\parindent}{1.5em}\noindent
A covector is a machine that eats a vector and returns a number; this evaluation, called
\emph{contraction}, is the one operation from which most of what follows is built. We start from it,
develop tensors, the exterior algebra and the geometric (Clifford) algebra together with their duality,
and then pass to manifolds, differential forms and the Lie derivative. Contraction with a velocity
defines the integral of a $1$-form; contraction with a tangent frame defines the integral of a
$k$-form; and the theorems of Stokes and Gauss are understood today as a single, metric-free
statement, ``$\dd$ is the adjoint of $\partial$'', of which the classical theorems are special cases.
Adding a connection gives covariant differentiation, geodesics and the Riemann curvature tensor,
which we follow to geodesic deviation and Einstein's equations. A physical part sorts quantities
of mechanics, thermodynamics, electromagnetism, relativity and fluids into those that are intrinsic
forms and those that genuinely depend on a metric; a complex-analysis part shows that the residue
theorem is Stokes' theorem for a closed complex $1$-form on a region with holes. Three \emph{synthesis
sections} and a fourth synthesis subsection, together with gold boxes throughout, compare the three languages that recur throughout --
index/tensor calculus, differential forms, and geometric algebra -- and reconcile their conventions.
A closing glossary lists the symbols and terms and collects the main identities with complete proofs.
\end{minipage}
\end{center}
